Este capítulo delimita con precisión qué comprende el presente trabajo de graduación, estableciendo sus límites funcionales, técnicos y operativos, así como las condiciones, restricciones y supuestos bajo los cuales se desarrolla. La delimitación aquí definida es consistente con los objetivos planteados y sirve como referencia directa para evaluar, al cierre del proyecto, si lo entregado corresponde a lo comprometido.

\section{Alcance funcional}

logSguarDian es una librería npm que implementa un middleware de tipo \textit{Runtime Application Self-Protection} (RASP) para aplicaciones Node.js construidas sobre el framework Express. Su alcance funcional comprende cuatro elementos concretos. Primero, la intercepción de cada petición HTTP entrante antes de que alcance la lógica de negocio de la aplicación anfitriona. Segundo, la extracción y vectorización de un conjunto de características numéricas del objeto \texttt{req} de Express, organizadas en seis grupos funcionales orientados a capturar los marcadores semánticos de las cuatro categorías de ataque cubiertas, además de características estadísticas transversales y metadatos de la petición. Tercero, la inferencia sobre un modelo híbrido compuesto por un Random Forest (componente supervisado) y un Isolation Forest (componente no supervisado), ambos serializados en formato ONNX y ejecutados de forma asíncrona mediante \texttt{worker\_threads} sin bloquear el \textit{Event Loop} del proceso anfitrión. Cuarto, la clasificación de cada petición como tráfico legítimo o como una de las cuatro categorías de ataque cubiertas, con la generación correspondiente de una respuesta de bloqueo configurable o de una alerta registrada.

El alcance funcional incluye además una interfaz de línea de comandos con cuatro grupos de subcomandos (\texttt{config}, \texttt{endpoints}, \texttt{attacks}, \texttt{webhooks}) que permiten al desarrollador anfitrión obtener visibilidad sobre el tráfico clasificado, perfilar rutas con mayor densidad de ataques e integrar la librería con sistemas externos de notificación mediante webhooks, sin requerir modificaciones a la arquitectura de la aplicación protegida.

\section{Alcance técnico}

El componente de Machine Learning documentado en este trabajo comprende: el diseño e implementación del pipeline de ingeniería de características que transforma el objeto \texttt{req} de Express en un vector numérico de dimensión fija; la construcción de un corpus de entrenamiento compuesto predominantemente por datasets públicos curados (entre ellos SR-BH~2020/CAPEC, modsec-learn, registros de honeypot de ModSecurity, SecLists, PayloadsAllTheThings y conjuntos de Kaggle, cada uno con licencia y procedencia verificadas), complementado con un conjunto acotado de tráfico sintético generado en un entorno de prueba controlado (Docker Compose) para cubrir vacíos puntuales de cobertura detectados durante la validación del modelo, con etiquetado determinista basado en el conocimiento del proceso de generación para esa porción sintética; el entrenamiento, ajuste de hiperparámetros y evaluación del modelo híbrido Random Forest e Isolation Forest sobre dicho corpus; la serialización del modelo en formato ONNX mediante \texttt{sklearn2onnx}, con verificación de paridad de predicción respecto al modelo original en Python; y la integración del modelo serializado en la librería npm mediante el runtime \texttt{onnxruntime-node}, con inferencia ejecutada en un \textit{worker pool} dedicado.

La librería está delimitada técnicamente a aplicaciones que utilicen Node.js en versión igual o superior a 12 (requisito mínimo de \texttt{worker\_threads} de forma estable) junto con el framework Express como capa de enrutamiento HTTP. El entrenamiento de los modelos se realiza íntegramente en Python sobre hardware Apple Silicon local, sin dependencia de infraestructura de cómputo en la nube.

\section{Alcance operativo y de evaluación}

La evaluación del sistema se realiza en un entorno de prueba controlado y reproducible, compuesto por una aplicación Express deliberadamente vulnerable desplegada mediante Docker Compose. Sobre este entorno se definieron tres configuraciones de evaluación progresivas: una configuración base sin ningún mecanismo de protección, que confirma la explotabilidad de la aplicación de referencia; una configuración con logSguarDian activo como única capa de protección; y una configuración de defensa en capas que combina un Web Application Firewall perimetral (Nginx con ModSecurity y el Core Rule Set de OWASP) junto con logSguarDian operando como segunda capa. Esta tercera configuración delimita explícitamente el alcance operativo del proyecto a un escenario de defensa en profundidad, evaluando la capacidad de detección de logSguarDian de forma independiente sobre el tráfico que logra evadir o queda fuera de la cobertura de un WAF basado en firmas, en lugar de plantear la evaluación como una comparación de qué componente es individualmente superior.

El alcance de esta evaluación se limita al entorno de prueba controlado descrito; no se contempla dentro del alcance del proyecto un despliegue en un entorno de producción real de terceros como criterio de aceptación del sistema. Cualquier validación adicional con organizaciones externas se documenta como evidencia contextual complementaria y no como parte de los criterios de aceptación verificables del proyecto.

\section{Elementos fuera del alcance}

La Tabla \ref{cuadro:fuera_alcance} enumera las exclusiones explícitas del proyecto, cada una acompañada de su justificación técnica, con el fin de delimitar con precisión la frontera entre lo que logSguarDian cubre y lo que queda deliberadamente fuera de su diseño.

\begin{table}[h]
\small
\begin{tabular}{|p{3.6cm}|p{9.4cm}|}
\hline
\textbf{Ámbito excluido} & \textbf{Razón técnica} \\ \hline
Denegación de servicio (DDoS) y fuerza bruta & Ataques inter-request: su naturaleza maliciosa reside en patrones de frecuencia y volumen acumulados entre múltiples peticiones, incompatibles con un vector de características que opera petición a petición. \\ \hline
BOLA, Mass Assignment, SSRF, deserialización insegura & Requieren conocimiento del modelo de datos específico de la aplicación anfitriona; no generan patrones semánticos detectables sin contexto de la lógica de negocio. \\ \hline
Inyecciones NoSQL (MongoDB u otras) & Sintaxis de payload estructuralmente distinta a las inyecciones SQL cubiertas; requeriría un corpus de entrenamiento y marcadores de características independientes. \\ \hline
Modelos de aprendizaje profundo de alta capacidad (CNN, RNN, Transformers) & Costo de inferencia incompatible con el presupuesto de latencia del middleware; tamaño de modelo serializado incompatible con las convenciones de un paquete npm de uso general. \\ \hline
Aprendizaje continuo y reentrenamiento automático embebido en el proceso de la aplicación anfitriona & Requeriría infraestructura de almacenamiento de muestras y validación previa al despliegue dentro del propio proceso del middleware, incompatible con el principio de operación autónoma sin dependencias externas de la librería en tiempo de ejecución. \\ \hline
Seguridad de red (Capas 3 y 4 del modelo OSI) & logSguarDian opera exclusivamente en la Capa 7, dentro del proceso Node.js; la mitigación volumétrica de red es responsabilidad de controles perimetrales complementarios. \\ \hline
Funcionalidades administrativas de WAF empresarial (gestión centralizada de reglas, dashboards, integración con CDN) & Requieren infraestructura de servicio central, incompatible con el principio de operación sin dependencias externas que fundamenta la propuesta de valor de la librería. \\ \hline
Compatibilidad con runtimes distintos de Node.js/Express & El mecanismo de integración es específico al sistema de middleware de Express; soporte para otros lenguajes o frameworks constituiría un entregable independiente. \\ \hline
\end{tabular}
\caption[Elementos fuera del alcance del proyecto]{Elementos explícitamente fuera del alcance del proyecto, con su justificación técnica. Elaboración propia.}
\label{cuadro:fuera_alcance}
\end{table}

\section{Restricciones, condiciones y supuestos}

El desarrollo del proyecto está sujeto a las siguientes restricciones y supuestos, que delimitan el margen de diseño y ejecución disponible. En cuanto a restricciones de tiempo, el proyecto se ejecuta dentro del período comprendido entre mayo y octubre de 2026, correspondiente al calendario del trabajo de graduación de la Facultad de Ingeniería de la Universidad del Valle de Guatemala; cualquier línea de trabajo no completable dentro de dicho período se documenta explícitamente como trabajo futuro y no como un entregable pendiente. En cuanto a restricciones de recursos, el entrenamiento de los modelos se realiza sobre hardware local sin presupuesto asignado a instancias de cómputo en la nube, condición verificada como viable dado el tamaño de los modelos de ensamble seleccionados (Random Forest e Isolation Forest) frente a alternativas de mayor capacidad computacional. En cuanto a restricciones de dependencias externas, la librería está diseñada para operar sin comunicación con servicios de inferencia remotos ni infraestructura de terceros durante su ejecución en producción, lo cual delimita el diseño hacia modelos de inferencia local suficientemente ligeros para residir en el mismo proceso que la aplicación anfitriona.

Se asume, además, que el conjunto de herramientas de simulación de ataques y generación de tráfico legítimo empleadas para construir el corpus de entrenamiento produce muestras suficientemente representativas de las cuatro categorías de ataque cubiertas; esta condición fue verificada mediante el análisis exploratorio del corpus resultante antes de proceder a la fase de entrenamiento. Se asume también que la arquitectura de middleware de Express constituye un punto de integración representativo del patrón dominante de desarrollo de APIs REST en el ecosistema de startups tecnológicas guatemaltecas descrito en la Justificación de este trabajo.

\section{Coherencia y viabilidad del alcance}

El alcance definido en este capítulo mantiene coherencia directa con los objetivos específicos del proyecto: el alcance funcional y técnico corresponde a los objetivos de construcción del pipeline de características y del modelo híbrido, mientras que el alcance operativo de evaluación en tres configuraciones progresivas corresponde al objetivo de evaluación del desempeño integral del sistema. Las exclusiones documentadas en la Tabla \ref{cuadro:fuera_alcance} son decisiones de diseño deliberadas, tomadas con anterioridad al inicio del desarrollo y orientadas a mantener un alcance verificable dentro del tiempo y los recursos disponibles, no limitaciones que el proyecto haya dejado sin resolver. La viabilidad de este alcance está respaldada por la naturaleza acotada de las cuatro categorías de ataque seleccionadas --todas de naturaleza intra-request, evaluables mediante el mismo vector de características sin requerir estado persistente entre peticiones-- y por la selección de algoritmos de complejidad computacional lineal, cuya factibilidad de entrenamiento e inferencia bajo las restricciones de hardware y tiempo del proyecto fue confirmada antes de comprometer dicho alcance como parte de los objetivos específicos del trabajo.
